% !TeX root = ../教材题库_全部条目_转写工作稿.tex
% Source: XJ-B1-C04-HW-P148-Y1232
% 第4章 / 4.5 函数模型及其应用 / 教材 p.148、149 / PDF 第155、156页
% OCR未校对：公式与文字均可能有误；未完成数学TeX转写；禁止直接发给学生
% 原题及相邻原解截图（核对用）：../题目原图/XJ-B1-C04-HW-P148-Y1232-01.webp, ../题目原图/XJ-B1-C04-HW-P148-Y1232-02.webp
% 本文件只在转写完成后作为题干使用；OCR 条目尚不是可靠题干。
\begin{bankitem}{习题4.5 第6题}{第4章 / 4.5 函数模型及其应用 / 教材 p.148、149 / PDF 第155、156页}{XJ-B1-C04-HW-P148-Y1232}{OCR未校对：公式与文字均可能有误；未完成数学TeX转写；禁止直接发给学生}
\begin{ocrdraft}
[教材p.148 / PDF第155页]
\par
6．某公司为了实现1000万元利润的目标，准备制定一个激励销售人员的奖励
\par
方案：在销售利润达到10万元时，按销售利润进行奖励，且奖金万元）随销售利
\par
润云万元）的增加而增加，但奖金总数不超过5万元，同时奖金不超过利润的25％．
\par
现有以下三个奖励模型：
\par
（1）3，一0．25；
\par
（2）3，一log7工十1；
\par
（3）3，一1.002飞
\par
其中哪个模型能符合公司的要求（logl.“5805．52，10g过000
\par
3．55）？
\par
[教材p.149 / PDF第156页]
\par
温故而知新
\end{ocrdraft}
\end{bankitem}
